\documentclass[conference]{IEEEtran}

\usepackage[T1]{fontenc}
\usepackage[utf8]{inputenc}
\usepackage{graphicx}
\usepackage{booktabs}
\usepackage{hyperref}
\usepackage{xcolor}

\newcommand{\todo}[1]{\textcolor{red}{[#1]}}

\title{Adversarial Vulnerabilities in LLM-Based Chart Reading}

\author{
  \IEEEauthorblockN{\todo{Author}}
  \IEEEauthorblockA{\todo{Institution}}
}

\begin{document}
\maketitle

\begin{abstract}
LLM-based agents are increasingly asked to read data visualizations,
yet the presentation layer they receive is controlled by a third party.
We present \textit{vis-attack}, a corpus of 50 adversarially manipulated
charts across four charting libraries (D3, Plotly, Chart.js, Vega-Lite),
evaluated against six local open-weight models under four perception
conditions: raw HTML source, rendered screenshot, both combined, and a
two-agent pipeline. Tooltip/hover-text manipulation achieves 100\%
attack success rate across all libraries and conditions. We additionally
report a clean-chart capability baseline covering all ten Amar et al.\
low-level analytic tasks~\cite{amar2005}. Our main finding is that
models uniformly trust explicitly labeled values over structural data,
regardless of how the chart is delivered.
\end{abstract}

% ---------------------------------------------------------------
\section{Motivation}

LLMs are deployed in pipelines that read visualizations: BI copilots,
browsing agents, accessibility tools. In all these settings the
model's perception is mediated by source or rendered pixels that a
third party controls. If that presentation layer can be manipulated
to misrepresent the underlying data, the model produces a confident
wrong answer. We empirically measure how often this succeeds, across
libraries, models, and modalities.

% ---------------------------------------------------------------
\section{Research Questions}

\begin{itemize}
  \item \textbf{RQ1.} How often do adversarial chart manipulations
        fool LLM agents, and does the rate depend on perception
        modality?
  \item \textbf{RQ2.} Which attack types transfer across libraries
        and modalities?
  \item \textbf{RQ3.} Does giving a model more information
        (source + screenshot) improve robustness?
  \item \textbf{RQ4.} Does splitting perception and reasoning into
        two agents add a useful cross-check?
\end{itemize}

% ---------------------------------------------------------------
\section{Methods}

\subsection{Corpus}
Each library has one honest \texttt{clean.html} baseline and
11--14 attack variants. Every attack reuses the same underlying
dataset but manipulates one presentation channel (tooltip callback,
axis tick, hidden DOM layer, annotation, etc.) so a viewer relying
on that channel answers a fixed question incorrectly. Ground truth
is stored in a sidecar \texttt{.meta.json} never served to the model.

\subsection{Perception Conditions}

\begin{description}
  \item[\texttt{raw\_source}] Raw HTML/JS/SVG text, no rendering.
  \item[\texttt{vision\_screenshot}] Playwright screenshot only.
        Tooltip attacks are hovered before capture.
  \item[\texttt{multimodal\_combined}] Screenshot + source in one
        prompt.
  \item[\texttt{dual\_agent}] Vision model extracts facts from the
        screenshot; text model reasons over those facts alongside
        source-extracted facts; a third call reconciles both.
\end{description}

\subsection{Models}
All models served locally via Ollama on an NVIDIA V100 (SLURM).
Text: mistral:7b, llama3:8b, qwen2.5:7b.
Vision: llava:13b, qwen2.5vl.
Dual-agent: qwen2.5vl + qwen2.5:7b.

\subsection{Metric}
Attack Success Rate (ASR) = wrong rate on attack minus wrong rate
on clean baseline. Ambiguous responses are excluded. Numeric answers
graded within 1\%; categorical by substring match.

\subsection{Capability Baseline}
Separately, we run all ten Amar et al.\ tasks on each library's
clean chart under all four conditions, reporting raw accuracy.
Tasks are grouped into three tiers: Tier 1 (Retrieve Value, Find
Anomalies), Tier 2 (Find Extremum, Filter, Determine Range), Tier 3
(Sort, Compute Derived Value, Characterize Distribution, Cluster,
Correlate).

% ---------------------------------------------------------------
\section{Results}

\subsection{Attack Success Rate}

\begin{table}[h]
\caption{Mean ASR by library and condition.}
\label{tab:asr}
\centering
\small
\begin{tabular}{lrrr}
\toprule
Library & raw\_source & multimodal & dual\_agent \\
\midrule
D3        & $+46.7\%$ & $+35.7\%$ & $+42.9\%$ \\
Plotly    & $+37.1\%$ & $+15.0\%$ & $+38.9\%$ \\
Chart.js  & $0.0\%$   & $-18.2\%$ & $+27.3\%$ \\
Vega-Lite & $+2.8\%$  & $+12.5\%$ & $+25.0\%$ \\
\bottomrule
\end{tabular}
\end{table}

By model in \texttt{raw\_source}: llama3:8b $+28.2\%$,
qwen2.5:7b $+22.5\%$, mistral:7b $+2.5\%$.

Tooltip attacks (all four libraries' attack\_01) achieved
100\% ASR in every condition. Prompt-injection attacks via
chart title or metadata achieved 0\% ASR everywhere.

The \texttt{vision\_screenshot} condition cannot be interpreted for
robustness: llava:13b answered only $\sim$8.5\% of clean-baseline
questions correctly, leaving no headroom for attacks to worsen.

\subsection{Capability Baseline}

\begin{table}[h]
\caption{Mean accuracy by tier — D3 clean chart.}
\label{tab:cap}
\centering
\small
\begin{tabular}{llccc}
\toprule
Condition & Model & T1 & T2 & T3 \\
\midrule
raw\_source & mistral:7b        & 100\% & 0\%   & 60\% \\
            & llama3:8b         & 100\% & 67\%  & 40\% \\
            & qwen2.5:7b        & 100\% & 33\%  & 80\% \\
\midrule
vision      & llava:13b         & 0\%   & 33\%  & 40\% \\
            & qwen2.5vl         & 80\%  & 100\% & 80\% \\
\midrule
multimodal  & llava:13b         & 80\%  & 33\%  & 40\% \\
            & qwen2.5vl         & 100\% & 100\% & 80\% \\
\midrule
dual\_agent & qwen2.5vl+qwen2.5 & 100\% & 100\% & 80\% \\
\bottomrule
\end{tabular}
\end{table}

Compute Derived Value achieved 0\% for every model and condition,
consistent with Xu \& Wall's finding for GPT-4~\cite{xuwall2024}.
qwen2.5vl is the strongest vision model, matching raw-source
performance on Tier 2.

\todo{Insert figures: (1) example clean vs.\ attack chart pair;
(2) ASR heatmap across attacks; (3) capability bar chart by task.}

% ---------------------------------------------------------------
\section{Discussion}

\textbf{Tooltip trust is the dominant vulnerability.} Models prefer
explicitly labeled values over structural data. Xu \& Wall found
that \emph{presence} of labels helps LLMs read charts; we find that
\emph{falsified} labels are the single most reliable attack channel.
The two results are two sides of the same property.

\textbf{More information does not add robustness.} Multimodal
prompts still fell for tooltip attacks at 100\% ASR. A model given
contradictory signals did not flag the discrepancy.

\textbf{Dual-agent is the most susceptible condition.} Errors at
the vision-extraction stage propagate as unquestioned facts to the
reasoning stage. The cross-check only works when the extraction
model flags uncertainty — which it does not when fooled.

\textbf{Model choice matters more than condition.} mistral:7b's
raw\_source ASR (+2.5\%) is roughly one-tenth of llama3:8b's
(+28.2\%). Mitigation evaluations must control for model.

% ---------------------------------------------------------------
\section{Sources of Error}

\begin{itemize}
  \item \textbf{Negative ASR.} 15 attacks produced negative ASR,
        where the attack page was easier to answer than the clean
        baseline. These reflect ambiguous clean-baseline designs,
        not model robustness.
  \item \textbf{Single chart and question per attack.} Results may
        not generalize to other datasets or chart types.
  \item \textbf{Capability baseline: $n{=}1$.} One instantiation
        per task; 0\%/100\% results are suggestive, not robust.
  \item \textbf{Infrastructure exclusions.} llama3.2-vision
        returned HTTP 500 for all multimodal trials (excluded).
  \item \textbf{Local models only.} Frontier models (GPT-4o,
        Claude) may behave differently.
  \item \textbf{Zero-shot only.} Chain-of-thought or tool-augmented
        prompting could change both capability and susceptibility.
\end{itemize}

% ---------------------------------------------------------------
\section{Conclusion}

Tooltip manipulation achieves 100\% ASR across all four libraries
and every interpretable condition. Neither multimodal fusion nor
agent decomposition provides meaningful mitigation. Any LLM pipeline
that reads charts must treat labeled values as untrusted input.

\bibliographystyle{IEEEtran}
\begin{thebibliography}{9}

\bibitem{amar2005}
R.~Amar, J.~Eagan, J.~Stasko,
``Low-level components of analytic activity in information visualization,''
\emph{IEEE InfoVis}, 2005.

\bibitem{xuwall2024}
Z.~Xu and E.~Wall,
``Exploring the capability of LLMs in performing low-level visual analytic
tasks on SVG data visualizations,''
\emph{IEEE VIS (Short)}, 2024.

\end{thebibliography}

\end{document}
